% Generated by scripts/assemble.py; corrections belong in source/corrections.json.
\examhead{1}{北京师范大学 2024 年数学分析试卷}{9}
% p010-b01; source PDF page 10
\blockanchor{p010-b01}
\question{1}设  $a_{1} > 0$ ,  $a_{n+1} = \ln(1 + a_{n}) (n = 1, 2, \cdots)$ , 求  $\lim_{n \to \infty} a_{n}$  及  $\lim_{n \to \infty} na_{n}$ .\par
% p010-b02; source PDF page 10
\blockanchor{p010-b02}
\question{2}设 $f(x)$ 在 $[a, +\infty)$ 上连续, 且当 $x \to +\infty$ 时有渐近线 $y = bx + c$, 证明: $f(x)$ 在 $[a, +\infty)$ 上一致连续.\par
% p010-b03; source PDF page 10
\blockanchor{p010-b03}
\question{3}证明: 方程 $3^{x} + 7^{x} = 2 \cdot 5^{x}$ 的实根只有 0 和 1.\par
% p010-b04; source PDF page 10
\par\Needspace{6\baselineskip}
\blockanchor{p010-b04}
\question{4}设 $f(x)$ 在 $[0,1]$ 上连续且恒正, 证明: 对任意的正整数 $n$, 存在 $\xi_n \in (0,1)$, 使得\par
\nopagebreak[4]
% p010-b05; source PDF page 10
\blockanchor{p010-b05}
\displayformula{\frac {1}{n} \int_ {0} ^ {1} f (x) \mathrm{d} x = \int_ {0} ^ {\xi_ {n}} f (x) \mathrm{d} x + \int_ {1 - \xi_ {n}} ^ {1} f (x) \mathrm{d} x.}
% p010-b06; source PDF page 10
\blockanchor{p010-b06}
且 $\lim_{n\to \infty}n\xi_n = \frac{\int_0^1f(x)\mathrm{d}x}{f(0) + f(1)}.$\par
% p010-b07; source PDF page 10
\blockanchor{p010-b07}
\question{5}设 $f(x)$ 是以 $2\pi$ 为周期的连续函数, 记 $f_{n}(x) = n \int_{x}^{x + \frac{1}{n}} f(t) \mathrm{d}t$, 证明: $\{f_{n}(x)\}$ 为连续可微的 $2\pi$ 为周期的函数列, 且在 $\mathbb{R}$ 上一致收敛于 $f(x)$.\par
% p010-b08; source PDF page 10
\blockanchor{p010-b08}
\question{6}判断级数 $\sum_{n=2}^{\infty} \ln \left(1 + \frac{(-1)^n}{n^p}\right)$ 的敛散性, $p > 0$.\par
% p010-b09; source PDF page 10
\blockanchor{p010-b09}
\question{7}计算极限 $\lim_{n\to \infty}\int_0^1\frac{\sin^2nx}{1 + x^2}\mathrm{d}x.$\par
% p010-b10; source PDF page 10
\blockanchor{p010-b10}
\question{8}设 $f(x, y, z)$ 有连续的偏导数，且满足 $-y\frac{\partial f}{\partial x} + x\frac{\partial f}{\partial y} + b\frac{\partial f}{\partial z} \geq c$，其中 $b, c$ 为正实数，证明：当 $(x, y, z)$ 沿着 $x = b\cos t, y = b\sin t, z = bt, t \in [0, +\infty)$ 趋于无穷时，$f(x, y, z) \to +\infty$.\par
% p010-b11; source PDF page 10
\par\Needspace{6\baselineskip}
\blockanchor{p010-b11}
\question{9}设 $\Omega$ 是 $z = x^{2} + y^{2}$ 介于 $z = 0$ 和 $z = 2$ 之间的曲面, 求曲面积分\par
\nopagebreak[4]
% p010-b12; source PDF page 10
\blockanchor{p010-b12}
\displayformula{\iint_ {\Sigma} \left(x ^ {2} \cos \alpha + y ^ {2} \cos \beta + z ^ {2} \cos \gamma\right) d S.}
% p010-b13; source PDF page 10
\blockanchor{p010-b13}
其中 $(\cos\alpha,\cos\beta,\cos\gamma)$ 为 $\Sigma$ 的单位外法向量.\par
% p010-b14; source PDF page 10
\par\Needspace{6\baselineskip}
\blockanchor{p010-b14}
\question{10}计算曲线积分\par
\nopagebreak[4]
% p010-b15; source PDF page 10
\blockanchor{p010-b15}
\displayformula{\oint_ {L} \left(y ^ {2} - z ^ {2}\right) \mathrm{d} x + \left(z ^ {2} - x ^ {2}\right) \mathrm{d} y + \left(x ^ {2} - y ^ {2}\right) \mathrm{d} z.}
% p010-b16; source PDF page 10
\blockanchor{p010-b16}
其中 $L$ 为 $x^{2} + y^{2} + z^{2} = 4$ 与 $x^{2} + y^{2} = 4$ 的交线, 从 $z$ 轴正向看为逆时针方向. (可能有误)\par
